\documentclass[preprint,12pt]{elsarticle}
\pdfminorversion=7

\usepackage{booktabs}
\usepackage{array}
\usepackage{tabularx}
\usepackage{graphicx}
\usepackage{xcolor}
\usepackage{amsmath}
\usepackage[hidelinks]{hyperref}
\setlength{\emergencystretch}{2em}

\newcommand{\ownedfigure}[2]{%
\includegraphics[width=0.94\linewidth]{#2}}

\journal{Computers in Biology and Medicine}

\begin{document}

\begin{frontmatter}

\title{Reliability Validation of COVID-19 Respiratory-Audio Screening Under Temporal Drift, External Dataset Shift, and Metadata Confounding}

\author[inst1]{Author names withheld for review}
\address[inst1]{Institutional affiliation withheld for review}

\begin{abstract}
Respiratory audio is attractive for scalable COVID-19 screening, but a biomedical signal-processing model must be evaluated for reliability rather than only optimized on an internal benchmark. We built a modality-scoped validation framework using Coswara cough, breath, speech, and metadata as the source-domain dataset, and COUGHVID only as an independent cough-only transfer target. The pipeline audited label availability and signal quality, extracted openSMILE ComParE+IS10 acoustic descriptors, evaluated classical and boosted model families with source-domain feature selection, and compared participant-disjoint, time-aware, strict early-to-late temporal, external-transfer, and metadata-only protocols. The selected Coswara participant-split fusion model reached AUROC 0.897, AUPRC 0.863, balanced accuracy 0.825, and F1 0.752. Performance decreased in the time-stratified participant split (AUROC 0.849) and in strict early-to-late temporal validation (AUROC 0.698). In model-matched cough-only comparisons, Coswara internal AUROC was 0.849-0.868, whereas COUGHVID transfer AUROC was 0.523-0.543. Metadata-only controls were highly predictive, with full safe metadata AUROC 0.964, symptoms-only AUROC 0.932, and demographic/protocol-only AUROC 0.914. Inverse-propensity sensitivity reduced quality-weighted audio AUROC from 0.879 unweighted to 0.780-0.807 after measured adjustment. These results show that internally strong respiratory-audio discrimination can coexist with temporal instability, external dataset shift, and non-audio confounding. COVID-19 respiratory-audio studies should therefore report modality scope, temporal validation, external transfer, metadata controls, calibration behavior, and clinically defined operating points before making screening claims.
\end{abstract}

\begin{keyword}
COVID-19 \sep respiratory audio \sep cough analysis \sep temporal drift \sep confounding \sep external validation \sep biomedical signal processing
\end{keyword}

\end{frontmatter}

\section{Introduction}

Crowdsourced cough, breathing, and speech recordings created an appealing route to low-cost COVID-19 screening. A mobile recording can be collected without laboratory equipment, repeated over time, and analyzed with standard acoustic or deep-learning representations. The biomedical question, however, is not whether a classifier can separate positives and negatives in a convenient retrospective split. For screening, a model must preserve clinically meaningful behavior when the population, recording conditions, prevalence, symptom distribution, device mix, and calendar phase change.

COVID-19 respiratory audio is especially vulnerable to this gap between internal discrimination and screening reliability. Infection labels in observational datasets can be entangled with symptoms, testing availability, recruitment period, country, device, and recording protocol. Audio itself can also carry non-pathophysiological information: background noise, microphone characteristics, compliance with prompted maneuvers, or vocal effort may differ between groups. A signal-processing validation study must therefore ask whether a model is learning a portable respiratory signature, a source-dataset regularity, or a mixture of both.

This manuscript evaluates COVID-19 respiratory-audio modeling as a reliability problem. Coswara provides the source-domain setting for cough, breath, speech, temporal, and metadata analyses \citep{bhattacharya2023coswara}. COUGHVID is used as an independent external target only for cough transfer, because it does not provide matching breath and speech modalities \citep{orlandic2021coughvid}. The resulting claims are intentionally modality-scoped: Coswara supports multimodal source-domain and temporal validation, whereas COUGHVID supports cough-only external transfer.

The paper is therefore not organized around a new architecture alone. It is organized around the biomedical validation question: when a respiratory-acoustic classifier appears strong in one collection process, what evidence remains after signal-quality auditing, participant separation, temporal stress testing, external cough transfer, and non-audio metadata controls?

The study makes four contributions. First, it describes a quality-audited, participant-aware acoustic pipeline for crowdsourced COVID-19 respiratory audio. Second, it compares internal participant, time-stratified participant, strict early-to-late temporal, and external cough-transfer protocols using the same validation logic. Third, it reports metadata-only and temporal metadata ablation controls to quantify confounding risk. Fourth, it interprets lower temporal and external performance as scientifically meaningful evidence about screening reliability, rather than as a secondary negative result.

\section{Background and Related Work}

\subsection{Respiratory-Audio Datasets}

Several public or semi-public resources made COVID-19 audio screening studies possible. Coswara collected cough, breath, sustained-vowel, counting, symptom, and demographic/protocol information for remote SARS-CoV-2 screening research \citep{bhattacharya2023coswara}. COUGHVID focused on crowdsourced cough recordings with expert and metadata-derived labels, and is therefore well suited to cough-algorithm development but not to multimodal fusion claims \citep{orlandic2021coughvid}. COVID-19 Sounds and the associated Sounds of COVID-19 analyses broadened the field by collecting large-scale cough, breath, voice, and participant metadata under realistic digital-health conditions \citep{xia2021covid,han2022sounds}. These datasets are valuable precisely because they reflect scalable collection; the same scale also introduces label noise, protocol heterogeneity, and confounding that must be tested rather than assumed away.

\subsection{Handcrafted and Classical Acoustic Modeling}

Early COVID-19 cough studies showed that conventional acoustic features and machine-learning classifiers could produce high internal performance under specific collection and validation designs. Pahar et al. reported COVID-19 cough classification from global smartphone recordings in \textit{Computers in Biology and Medicine} \citep{pahar2021global}, and Chowdhury et al. used ensemble and multi-criteria decision-making approaches for cough-based detection \citep{chowdhury2022mcdm}. These studies motivate a careful handcrafted-feature baseline rather than a reflexive assumption that larger neural architectures are required. In the present work, openSMILE-style ComParE+IS10 descriptors provide a broad acoustic representation spanning spectral, cepstral, energy, voicing, and functional summary information \citep{eyben2010opensmile}.

\subsection{Deep, Transformer, and Multimodal Methods}

Deep transfer-learning and bottleneck-feature studies extended COVID-19 audio modeling from cough to cough, breath, and speech, including CBM work on multimodal respiratory audio \citep{pahar2022transfer}. Hierarchical spectrogram transformer models further show how architecture design can exploit local-to-global time-frequency structure in respiratory sounds \citep{aytekin2024hst}. These architecture-forward studies are important comparators, but they answer a different primary question from the present manuscript. Here, the central contribution is not a new neural architecture; it is the validation of acoustic models under participant, temporal, external, and metadata-confounding stress.

\subsection{Cross-Dataset, Temporal, and Confounding Evidence}

Recent work has made the reliability problem more explicit. Islam et al. studied robust cough-based COVID-19 detection across multiple datasets with deep neural decision tree and forest methods \citep{islam2026eswa}. Ganitidis et al. examined drift-adaptive COVID cough modeling under chronological splits and adaptation mechanisms \citep{ganitidis2025jmir}. Coppock et al. showed that audio-based classifiers may not improve realistic screening beyond simple symptom checkers when confounding is carefully evaluated \citep{coppock2024nmi}. Together, these studies shift the field from headline internal accuracy toward source shift, temporal drift, calibration, and clinical operating-point interpretation.

Table~\ref{tab:literature_context} summarizes the protocol context for the most relevant evidence. The table is not a leaderboard. Internal cross-validation, participant-disjoint validation, temporal validation, and external transfer answer different biomedical questions and should not be ranked as if they were interchangeable.

\begin{table}[t]
\centering
\caption{Protocol-oriented literature context for COVID respiratory-audio screening. The comparison separates dataset resources, feature-engineered models, deep and transformer models, cross-dataset studies, temporal-drift studies, and confounding analyses. Metrics from these papers are not treated as directly interchangeable unless the dataset, validation unit, modality, and endpoint match.}
\label{tab:literature_context}
\scriptsize
\setlength{\tabcolsep}{2.5pt}
\begin{tabularx}{\linewidth}{p{0.18\linewidth} p{0.22\linewidth} p{0.23\linewidth} X}
\toprule
Study & Dataset/modality & Main validation focus & Relevance to this manuscript \\
\midrule
Coswara \citep{bhattacharya2023coswara} & Cough, breath, speech, symptoms, metadata & Public source dataset & Source for multimodal acoustic modeling, temporal analysis, and metadata confounding controls. \\
COUGHVID \citep{orlandic2021coughvid} & Cough & Public cough dataset & Independent cough-only transfer target; not evidence for multimodal fusion. \\
COVID-19 Sounds \citep{xia2021covid} & Cough, breath, voice, metadata & Large-scale digital screening resource & Field context for scalable respiratory-audio collection. \\
Sounds of COVID-19 \citep{han2022sounds} & COVID-19 Sounds & Realistic audio-based digital testing & Motivates cautious interpretation under practical collection conditions. \\
Pahar et al. \citep{pahar2021global} & Smartphone cough & Internal cough classification & Demonstrates high internal cough performance under defined collection settings. \\
Pahar et al. \citep{pahar2022transfer} & Cough, breath, speech & Transfer learning and bottleneck features & Multimodal respiratory-audio precedent; our emphasis is validation under shift. \\
Chowdhury et al. \citep{chowdhury2022mcdm} & Cough & Ensemble machine learning & Supports multi-model comparison and thresholded metric reporting. \\
HST \citep{aytekin2024hst} & Respiratory sounds & Spectrogram transformer architecture & Architecture-forward comparator for learned time-frequency modeling. \\
DNDT/DNDF \citep{islam2026eswa} & Multiple cough datasets & Cross-dataset cough evaluation & Context for interpreting weak Coswara-to-COUGHVID transfer. \\
Drift framework \citep{ganitidis2025jmir} & COVID cough audio & Temporal drift and adaptation & Motivates strict early-to-late validation. \\
Confounding critique \citep{coppock2024nmi} & Vocal audio and symptoms & Symptom and confounding analysis & Motivates symptoms-only and demographic/protocol metadata controls. \\
\bottomrule
\end{tabularx}
\end{table}

\section{Datasets and Quality Control}

\subsection{Modality Scope}

Coswara was the source dataset. It provides participant-level labels and multiple respiratory or speech-style recording prompts, enabling source-domain multimodal analysis. Cough, breath, and speech are therefore evaluated together only for Coswara endpoints. COUGHVID was the external target. Because COUGHVID is cough-only, it was used only to test whether a Coswara-trained cough representation transfers to an independently collected cough corpus.

\begin{figure}[t]
\centering
\includegraphics[width=0.94\linewidth]{../common_figures/fig_dataset_quality_flow.pdf}
\caption{Dataset and quality-control flow for the source and external cohorts. Coswara recordings with known positive or negative labels support supervised cough, breath, speech, temporal, and metadata analyses. Unknown-label recordings are excluded from supervised metrics but remain visible in the descriptive audit. COUGHVID enters only as an external cough-only test set, so it cannot validate Coswara multimodal fusion.}
\label{fig:dataset_quality_flow}
\end{figure}

\begin{table}[t]
\centering
\caption{Coswara supervised label counts used in the dataset audit. Counts are separated by modality because speech contributes multiple prompted recordings per participant, whereas cough and breath contribute fewer recordings per participant. Unknown-label rows are descriptive and are not used in supervised performance estimates.}
\label{tab:dataset_counts}
\begin{tabular}{llrrr}
\toprule
Modality & Label & Recordings & Participants & Dataset \\
\midrule
Breath & Negative & 2866 & 1433 & Coswara \\
Breath & Positive & 1362 & 681 & Coswara \\
Breath & Unknown & 1264 & 632 & Coswara \\
Cough & Negative & 2865 & 1433 & Coswara \\
Cough & Positive & 1362 & 681 & Coswara \\
Cough & Unknown & 1264 & 632 & Coswara \\
Speech & Negative & 7164 & 1433 & Coswara \\
Speech & Positive & 3405 & 681 & Coswara \\
Speech & Unknown & 3160 & 632 & Coswara \\
\bottomrule
\end{tabular}
\end{table}

\subsection{Signal Quality Audit}

Crowdsourced respiratory recordings vary in duration, silence content, file integrity, and prompt adherence. Before supervised modeling, recordings were audited for basic usability. The audit separated usable recordings from mostly silent, corrupt, and very short files (Table~\ref{tab:quality}). These checks are not a substitute for clinical-grade acquisition, but they reduce the chance that the classifier is dominated by trivial file-quality artifacts.

\begin{table}[t]
\centering
\caption{Audio quality audit summary. The median duration and silence ratio describe the audio files assigned to each quality flag. The supervised pipeline emphasizes usable recordings while preserving the audit as evidence of crowdsourced signal variability.}
\label{tab:quality}
\begin{tabular}{lrrr}
\toprule
Quality flag & Recordings & Median duration (s) & Median silence ratio \\
\midrule
OK & 23539 & 8.789 & 0.106 \\
Mostly silence & 597 & 7.680 & 0.802 \\
Corrupt & 390 & 0.000 & 1.000 \\
Short & 186 & 0.597 & 0.000 \\
\bottomrule
\end{tabular}
\end{table}

\section{Methods}

\subsection{Acoustic Representation}

The final acoustic representation used openSMILE-style ComParE+IS10 descriptors \citep{eyben2010opensmile}. This choice reflects a signal-processing objective: capture a broad set of interpretable low-level descriptors and summary functionals from variable-length respiratory recordings, without requiring a large end-to-end neural model. Candidate representations and branches included spectrogram CNN/BiGRU models and WavLM self-supervised embeddings, but the final selected validation rows used the ComParE+IS10 top-800 feature strategy because it gave the strongest and most stable audited validation results. ComParE+IS10 was also used for the external-transfer endpoint because it is a standardized handcrafted representation aligned with much of the respiratory-audio baseline literature and can be applied to COUGHVID without target-domain tuning.

Feature selection was performed inside the modeling workflow and not on the external COUGHVID labels. The selected top-800 feature set was based on the Coswara source-domain training workflow and then frozen for downstream validation. This distinction matters because external transfer is intended to test portability, not to retune the model to the target dataset. For strict chronological validation, the top-800 representation should be interpreted as a frozen source-domain representation rather than a fully prospective feature-selection protocol re-estimated only inside the early-period cohort.

The WavLM and CNN-BiGRU rows are source-domain single-stream checks, not completed COUGHVID transfer experiments. They should therefore be read as model-development breadth rather than as external validation evidence. The stronger 0.897 AUROC endpoint is a validation-selected Coswara multimodal fusion row, whereas the WavLM and CNN rows are isolated modality branches.

\subsection{Model Bank and Fusion}

The model bank included boosted-tree, kernel, and regularized linear classifiers. The final source-domain rows used validation-selected fusion over modality scores. The existing participant split combined cough and speech with stacked logistic fusion. The time-stratified participant split combined cough, breath, and speech using uniform score averaging. The strict early-to-late temporal row selected a breath top-4 validation ensemble. These are reported as selected validation outcomes for their respective protocols, not as a claim that one fusion rule or modality set is universally optimal.

Fusion was always modality-scoped. Cough, breath, and speech fusion is evaluated only where all required Coswara modalities exist for the endpoint. COUGHVID external transfer evaluates cough models only. No result in this manuscript should be read as multimodal external validation on COUGHVID.

\begin{figure}[t]
\centering
\ownedfigure{End-to-end pipeline diagram}{../common_figures/fig_pipeline.pdf}
\caption{Signal-processing and validation pipeline. Raw Coswara audio and metadata pass through label and quality auditing, modality-specific acoustic feature extraction, model-bank evaluation, source-domain feature selection, optional source-domain fusion, and the reliability validation ladder. COUGHVID enters only at the external cough-transfer stage.}
\label{fig:pipeline}
\end{figure}

\begin{figure}[t]
\centering
\ownedfigure{Multimodal fusion schematic}{../common_figures/fig_multimodal_fusion.pdf}
\caption{Modality-scoped fusion design. Coswara endpoints can combine cough, breath, and speech scores because those modalities are available for the source cohort. COUGHVID contains cough recordings only, so it tests cough-domain transfer and cannot establish multimodal fusion generalization.}
\label{fig:fusion_schematic}
\end{figure}

\subsection{Validation Protocols}

The validation ladder was designed to expose increasingly difficult reliability threats.

\begin{enumerate}
\item The existing participant split tested source-domain discrimination with participant separation.
\item The time-stratified participant split preserved participant separation while balancing calendar structure.
\item The strict early-to-late temporal protocol trained on earlier participants and tested on later participants, directly probing temporal shift.
\item The Coswara-to-COUGHVID protocol trained on Coswara cough data and tested on independent COUGHVID cough recordings.
\item Metadata-only audits tested whether non-audio variables could predict infection labels without any acoustic signal.
\end{enumerate}

\begin{figure}[t]
\centering
\ownedfigure{Evaluation protocol ladder}{../common_figures/fig_evaluation_ladder.pdf}
\caption{Reliability validation ladder. Each step changes the interpretation of a COVID-19 respiratory-audio result: internal participant separation addresses participant leakage, time stratification addresses calendar imbalance, strict early-to-late validation addresses chronological shift, COUGHVID transfer addresses independent cough-domain shift, and metadata-only controls address shortcut and confounding risk.}
\label{fig:protocol_ladder}
\end{figure}

\subsection{Metrics}

AUROC measures rank discrimination over thresholds. AUPRC is reported because positive cases are clinically important, but it is prevalence-sensitive and should be interpreted with the target positive rate. Balanced accuracy averages sensitivity and specificity, making thresholded performance less dominated by the majority class. F1 summarizes precision and sensitivity at the selected threshold but does not include true negatives. Sensitivity and specificity are reported where available because screening interpretation depends on the chosen operating point. Calibration metrics were considered in the broader analysis workflow, but the central claims in this manuscript rely on the final result tables that are available for all validation endpoints.

\subsection{Claim Boundary for Biomedical Interpretation}

Each validation protocol supports a different biomedical interpretation (Table~\ref{tab:claim_boundary}). The participant split tests source-domain discrimination, the time-aware and strict temporal splits test chronological stability, COUGHVID tests cough-only external transfer, and metadata-only models test shortcut risk. This separation is important because the same acoustic representation can appear clinically promising under one protocol and unreliable under another.

\begin{table}[t]
\centering
\caption{Biomedical interpretation boundary for each validation protocol. The table separates source-domain signal detection from future-period robustness, external cough portability, and metadata-confounding risk.}
\label{tab:claim_boundary}
\scriptsize
\setlength{\tabcolsep}{2.5pt}
\begin{tabularx}{\linewidth}{p{0.22\linewidth} p{0.20\linewidth} p{0.17\linewidth} X}
\toprule
Protocol evidence & Biomedical question & Main result & Interpretation boundary \\
\midrule
Participant-disjoint Coswara fusion & Is there source-domain respiratory-audio signal without participant overlap? & AUROC 0.897 & Supports feasibility inside Coswara only. \\
Time-aware and strict temporal validation & Is the source-domain signal stable across calendar structure? & AUROC 0.849 and 0.698 & Indicates temporal degradation and incomplete future-period robustness. \\
COUGHVID external transfer & Does a Coswara cough branch transfer to an independent cough corpus? & AUROC 0.543 & Indicates weak external cough portability under direct transfer. \\
Metadata-only controls & Can non-audio variables predict labels without respiratory sound? & AUROC 0.964 & Indicates substantial confounding and shortcut-learning risk. \\
\bottomrule
\end{tabularx}
\end{table}

\section{Results}

\subsection{Internal Discrimination Declines Under Time-Aware Validation}

The selected Coswara participant-split model showed strong source-domain discrimination (Table~\ref{tab:final_validation}). The cough+speech stacked logistic fusion row reached AUROC 0.897, AUPRC 0.863, balanced accuracy 0.825, and F1 0.752. This result supports the presence of label-associated acoustic information in Coswara under participant-disjoint internal validation.

The same evidence becomes more cautious when calendar structure is introduced. The time-stratified participant split decreased to AUROC 0.849 and balanced accuracy 0.783. Under strict early-to-late temporal validation, AUROC decreased to 0.698 and balanced accuracy to 0.656. The strict temporal AUPRC remained high at 0.896 because the temporal test set had high positive prevalence; it should therefore be interpreted together with AUROC, balanced accuracy, sensitivity, specificity, and the validation design.

\begin{table}[t]
\centering
\caption{Final ComParE+IS10 validation results across participant, temporal, and external-transfer settings. Metrics are rounded to three decimals from the source result tables. Coswara rows evaluate source-domain multimodal or modality-specific endpoints; the COUGHVID row evaluates cough-only external transfer.}
\label{tab:final_validation}
\resizebox{\linewidth}{!}{%
\begin{tabular}{llrrrrrrr}
\toprule
Protocol & Model/fusion & AUROC & AUPRC & Bal. acc. & F1 & Sens. & Spec. & $n$ \\
\midrule
Existing participant split & Cough+speech stacked logistic fusion & 0.897 & 0.863 & 0.825 & 0.752 & 0.833 & 0.816 & 314 \\
Time-stratified participant split & Cough+breath+speech uniform mean & 0.849 & 0.783 & 0.783 & 0.705 & 0.710 & 0.857 & 431 \\
Strict early-to-late temporal & Breath top-4 validation ensemble & 0.698 & 0.896 & 0.656 & 0.751 & 0.646 & 0.667 & 411 \\
COUGHVID external transfer & Cough LightGBM SMOTE f80 & 0.543 & 0.040 & 0.510 & 0.058 & 0.084 & 0.936 & 8331 \\
\bottomrule
\end{tabular}}
\end{table}

\begin{figure}[t]
\centering
\ownedfigure{Temporal degradation plot}{../common_figures/fig_validation_degradation.pdf}
\caption{Validation degradation across reliability settings. Bars show AUROC for selected final rows: strong source-domain Coswara discrimination, lower time-aware and strict temporal performance, and near-random Coswara-to-COUGHVID cough transfer. These selected endpoints differ in modality and fusion rule, so the plot is interpreted as claim-severity evidence rather than a fixed-model ablation; Table~\ref{tab:fixed_cough_ladder} reports the matched cough-only model-family ladder.}
\label{fig:temporal_degradation}
\end{figure}

\subsection{Paper-Comparable Cross-Validation Is Context, Not Deployment Evidence}

For comparison with prior internal-validation studies, paper-comparable 10-fold cross-validation produced cough AUROC 0.819 $\pm$ 0.029 and AUPRC 0.728 $\pm$ 0.036 for an RBF support-vector classifier. A similar LightGBM cough row reached AUROC 0.819 $\pm$ 0.027 and AUPRC 0.732 $\pm$ 0.035. Speech rows were also moderate, with the best speech AUROC 0.806 $\pm$ 0.014 and AUPRC 0.705 $\pm$ 0.021. These rows are useful for literature context, but they are not the main evidence for clinical reliability because recording-level cross-validation does not test chronological or external transfer.

\subsection{External COUGHVID Transfer Is Weak and Clinically Informative}

Direct Coswara-to-COUGHVID transfer produced near-random discrimination in the cough-only setting (Table~\ref{tab:coughvid_transfer}). The best final ComParE+IS10 external row reached AUROC 0.543, AUPRC 0.040, balanced accuracy 0.510, and F1 0.058 over 8331 COUGHVID samples. The target positive prevalence in the support analysis was approximately 0.034, so the AUPRC is only modestly above prevalence. This endpoint deliberately tests a frozen handcrafted ComParE+IS10 pipeline rather than every possible neural representation; it is a standardized cough-transfer stress test for a common paralinguistic-feature paradigm under maximum dataset shift.

This lower external performance is scientifically meaningful. It shows that a model selected on Coswara cough acoustics does not automatically learn a portable COVID-19 cough signature. COUGHVID differs in collection pathway, label process, acoustic quality, class prevalence, and cough-only modality scope. A reliable screening model must survive exactly these kinds of changes or clearly state that its use is restricted to the source setting.

Because COUGHVID is cough-only, the fair comparison is not Coswara multimodal fusion against COUGHVID transfer. The selected validation ladder is also not a fixed-model ablation because modality and fusion can change across selected endpoints. Table~\ref{tab:cough_matched_transfer} therefore compares cough-only Coswara internal rows with cough-only COUGHVID transfer rows for the same model families. Coswara cough-only internal AUROC ranged from 0.849 to 0.868, while COUGHVID transfer AUROC ranged from 0.523 to 0.543. Table~\ref{tab:fixed_cough_ladder} extends the matched comparison across the full validation ladder while holding modality fixed to cough and model family fixed. The weak external result is therefore not explained only by missing breath and speech modalities; the matched cough branch itself does not transfer robustly.

\begin{table}[t]
\centering
\caption{COUGHVID cough-only external transfer. All rows use the ComParE+IS10 top-800 feature strategy selected from the Coswara source-domain training workflow and are evaluated on independent COUGHVID cough labels. High specificity with very low sensitivity, as in the LightGBM row, is not an acceptable screening tradeoff without a prospectively justified operating point.}
\label{tab:coughvid_transfer}
\resizebox{\linewidth}{!}{%
\begin{tabular}{lrrrrrrr}
\toprule
Model & AUROC & AUPRC & Bal. acc. & F1 & Sens. & Spec. & $n$ \\
\midrule
LightGBM SMOTE f80 & 0.543 & 0.040 & 0.510 & 0.058 & 0.084 & 0.936 & 8331 \\
XGBoost SMOTE f80 & 0.532 & 0.039 & 0.505 & 0.063 & 0.267 & 0.743 & 8331 \\
CatBoost SMOTE f80 & 0.531 & 0.040 & 0.506 & 0.059 & 0.144 & 0.869 & 8331 \\
SVC RBF f60 & 0.523 & 0.037 & 0.500 & 0.058 & 0.193 & 0.807 & 8331 \\
\bottomrule
\end{tabular}}
\end{table}

\begin{table}[t]
\centering
\caption{Model-matched cough-only comparison between Coswara internal validation and COUGHVID external transfer. Both sides use cough-only ComParE+IS10 top-800 representations, making this the modality-matched external-portability comparison.}
\label{tab:cough_matched_transfer}
\small
\begin{tabular}{lcccccc}
\toprule
Model & \multicolumn{3}{c}{Coswara cough-only} & \multicolumn{3}{c}{COUGHVID cough-only} \\
\cmidrule(lr){2-4}\cmidrule(lr){5-7}
 & AUROC & AUPRC & BA & AUROC & AUPRC & BA \\
\midrule
LightGBM SMOTE & 0.853 & 0.797 & 0.782 & 0.543 & 0.040 & 0.510 \\
SVC RBF & 0.868 & 0.812 & 0.795 & 0.523 & 0.037 & 0.500 \\
CatBoost SMOTE & 0.849 & 0.788 & 0.776 & 0.531 & 0.040 & 0.506 \\
XGBoost SMOTE & 0.850 & 0.780 & 0.776 & 0.532 & 0.039 & 0.505 \\
\bottomrule
\end{tabular}
\end{table}

\begin{table}[t]
\centering
\caption{Fixed cough-only model-family ladder across validation tiers. Each row keeps the modality fixed to cough and follows the same model family across protocols. Models are retrained within the corresponding source-domain protocol; values are AUROC.}
\label{tab:fixed_cough_ladder}
\small
\resizebox{\linewidth}{!}{%
\begin{tabular}{lcccc}
\toprule
Model & Existing & Time-stratified & Strict temporal & COUGHVID \\
\midrule
LightGBM SMOTE & 0.853 & 0.830 & 0.604 & 0.543 \\
SVC RBF & 0.868 & 0.802 & 0.561 & 0.523 \\
CatBoost SMOTE & 0.849 & 0.826 & 0.614 & 0.531 \\
XGBoost SMOTE & 0.850 & 0.828 & 0.599 & 0.532 \\
\bottomrule
\end{tabular}
}
\end{table}

\subsection{Metadata-Only Controls Expose Strong Confounding Risk}

Metadata-only classifiers were more predictive than the final acoustic endpoints in the source-domain setting (Table~\ref{tab:confounding}). Full safe metadata reached AUROC 0.964 and balanced accuracy 0.890. Symptoms alone reached AUROC 0.932, and demographic/protocol variables alone reached AUROC 0.914. These models are not proposed as clinical predictors. They are negative controls for interpretation: if labels can be recovered from non-audio variables, then acoustic models trained on the same observational process may encode correlated shortcuts.

The temporal metadata analysis strengthened this concern. Removing recording month improved strict temporal metadata AUROC by 0.247 relative to the full temporal metadata row, indicating that calendar variables carried harmful temporal associations. Support analyses also identified recording year, recording month, country, age, and duration among prominent demographic/protocol attribution drivers. These findings connect the observed temporal degradation to plausible non-audio sources of shift.

\subsection{IPW Sensitivity Qualifies the Residual Audio Signal}

Inverse-propensity weighting was used as a measured-confounding sensitivity analysis for the quality-weighted audio fusion prediction set (Table~\ref{tab:ipw}). The unweighted row reached AUROC 0.879 and AUPRC 0.832 on 318 held-out samples. With cap 2 and 95th-percentile truncation, AUROC remained 0.807 with effective sample size 238.7. With cap 10 and 99th-percentile truncation, AUROC was 0.780, AUPRC was 0.537, effective sample size was 130.4, and maximum absolute residual standardized mean difference decreased to 0.382. The result is not causal proof. It supports a narrower conclusion: measured adjustment reduces apparent performance but leaves an above-chance residual signal that should be interpreted cautiously.

\begin{table}[t]
\centering
\caption{Inverse-propensity weighting sensitivity analysis for quality-weighted audio fusion. ESS denotes effective sample size; SMD denotes standardized mean difference.}
\label{tab:ipw}
\resizebox{\linewidth}{!}{%
\begin{tabular}{lrrrrrr}
\toprule
Weighting setting & AUROC & AUPRC & Bal. acc. & ESS & Mean abs. SMD & Max abs. SMD \\
\midrule
Unweighted & 0.879 & 0.832 & 0.804 & 318.0 & 0.224 & 1.384 \\
IPW cap 2, q=0.95 & 0.807 & 0.624 & 0.742 & 238.7 & 0.149 & 0.724 \\
IPW cap 5, q=0.95 & 0.793 & 0.573 & 0.731 & 173.2 & 0.129 & 0.523 \\
IPW cap 10, q=0.99 & 0.780 & 0.537 & 0.721 & 130.4 & 0.118 & 0.382 \\
\bottomrule
\end{tabular}}
\end{table}

\begin{table}[t]
\centering
\caption{Metadata confounding audit. The metadata rows are non-audio controls trained on Coswara metadata fields. Their purpose is to quantify label predictability from symptoms, demographic variables, and recording-protocol variables, not to propose a deployable diagnostic model.}
\label{tab:confounding}
\resizebox{\linewidth}{!}{%
\begin{tabular}{lrrrrrrl}
\toprule
Audit model & AUROC & AUPRC & Bal. acc. & F1 & Sens. & Spec. & Interpretation \\
\midrule
Full safe metadata & 0.964 & 0.928 & 0.890 & 0.849 & 0.858 & 0.922 & Non-audio metadata strongly predicts labels. \\
Symptoms only & 0.932 & 0.898 & 0.912 & 0.876 & 0.893 & 0.930 & Symptom fields are strong label proxies. \\
Demographic/protocol only & 0.914 & 0.737 & 0.827 & 0.751 & 0.864 & 0.790 & Non-symptom context remains predictive. \\
Remove month, temporal metadata & 0.779 & -- & -- & -- & -- & -- & AUROC improves by 0.247 versus full temporal metadata. \\
\bottomrule
\end{tabular}}
\end{table}

\begin{figure}[t]
\centering
\ownedfigure{Confounding evidence diagram/table}{../common_figures/fig_metadata_confounding.pdf}
\caption{Metadata confounding evidence. Non-audio metadata alone predicts Coswara labels strongly, and temporal metadata ablation shows that recording month can harm chronological generalization. These controls limit causal interpretation of acoustic performance in the source dataset.}
\label{fig:confounding_evidence}
\end{figure}

\subsection{Non-Final Deep Branches Did Not Resolve the Reliability Problem}

Representative CNN/spectrogram and WavLM self-supervised branches did not displace the final ComParE+IS10 validation rows (Table~\ref{tab:negative_branches}). Their inclusion is important because the observed temporal and external weaknesses should not be dismissed as a single-model artifact. These rows are source-domain single-modality checks, not COUGHVID transfer experiments and not direct ablations of multimodal fusion. In this analysis, adding practical deep branches did not remove the need for validation under shift.

\begin{table}[t]
\centering
\caption{Representative non-final model branches. These rows document model-development breadth and support the interpretation that reliability limits were not solved simply by switching to spectrogram CNN/BiGRU or practical WavLM embeddings.}
\label{tab:negative_branches}
\resizebox{\linewidth}{!}{%
\begin{tabular}{lllrrrrrr}
\toprule
Branch & Model & Modality & AUROC & AUPRC & Bal. acc. & F1 & Sens. & Spec. \\
\midrule
CNN/spectrogram & CNN-BiGRU & Breath & 0.710 & 0.590 & 0.649 & 0.550 & 0.684 & 0.614 \\
CNN/spectrogram & CNN-BiGRU & Cough & 0.730 & 0.556 & 0.686 & 0.586 & 0.680 & 0.693 \\
CNN/spectrogram & CNN-BiGRU & Speech & 0.676 & 0.490 & 0.638 & 0.524 & 0.579 & 0.698 \\
WavLM SSL transformer & WavLM base plus & Shallow breath & 0.790 & 0.654 & 0.731 & 0.627 & 0.804 & 0.657 \\
WavLM SSL transformer & WavLM base plus & Heavy cough & 0.789 & 0.696 & 0.730 & 0.636 & 0.673 & 0.787 \\
WavLM SSL transformer & WavLM base plus & Shallow cough & 0.790 & 0.731 & 0.713 & 0.614 & 0.608 & 0.817 \\
\bottomrule
\end{tabular}}
\end{table}

\section{Discussion}

The central result is not a single AUROC value. It is the pattern across validation settings: strong internal Coswara discrimination, reduced time-aware and strict temporal performance, weak model-matched COUGHVID cough transfer, and strong metadata-only predictability. This pattern is consistent with a respiratory-audio benchmark in which disease-associated acoustic information exists but is entangled with temporal, demographic, symptom, and protocol structure.

For biomedical signal processing, this has two consequences. First, source-domain fusion can be useful for studying multimodal respiratory audio, but it should not be confused with external multimodal validation unless the external dataset contains comparable modalities. Second, low external transfer is not merely a failed model-selection result. It is evidence about portability. A screening system that performs well only inside one collection process is not yet a general clinical tool, even if the internal split is participant-disjoint.

The metadata findings also change the clinical interpretation of acoustic performance. Symptoms may be clinically meaningful covariates in a diagnostic pathway, but they are not acoustic biomarkers. Demographic and protocol variables can reflect recruitment, testing access, and collection behavior. When these fields predict labels strongly, respiratory-audio papers must distinguish between audio-specific evidence, multimodal clinical prediction, and dataset-specific shortcut learning.

Future respiratory-audio screening studies should therefore report validation as part of the method, not as a post hoc robustness check. Minimum evidence should include participant separation, time-aware validation, at least one external dataset where modality-matched transfer is possible, metadata-only controls, calibration assessment, and clinically chosen thresholds. Prospective cohorts with standardized recording instructions, device and site stratification, event-level cough and breath segmentation, and subgroup calibration would provide a stronger basis for screening claims than retrospective leaderboard metrics alone.

\begin{figure}[t]
\centering
\includegraphics[width=0.94\linewidth]{../common_figures/fig_clinical_reliability_interpretation.pdf}
\caption{Clinical reliability interpretation. The evidence package supports cautious source-domain respiratory-audio modeling, not deployment-ready diagnosis. Internal performance must be interpreted alongside temporal degradation, weak cough-only external transfer, metadata confounding, and the absence of prospective calibration.}
\label{fig:clinical_reliability}
\end{figure}

\section{Limitations and Future Work}

This is a retrospective analysis of crowdsourced datasets, not a prospective clinical trial. SARS-CoV-2 labels, symptoms, recording conditions, devices, prompts, and recruitment pathways are not standardized in the way required for a clinical screening study. COUGHVID is cough-only, so the external experiment tests Coswara-to-COUGHVID cough transfer and cannot validate Coswara cough, breath, and speech fusion. The study therefore supports source-domain multimodal analysis and cough-only external-transfer analysis, but not external multimodal deployment.

Several analytical limitations remain. The study does not establish causal COVID-specific acoustic biomarkers, and the quality audit does not replace expert event segmentation or clinical adjudication. Metadata controls and inverse-propensity sensitivity analyses show measured confounding risk, but they do not fully decompose causal pathways among symptoms, testing behavior, recruitment period, country, device, and disease status. The cap-10 IPW row improves measured balance while reducing effective sample size to 130.4, so the retained AUROC 0.780 is supportive but not definitive. Structural confidence intervals were not computed for every final selected ComParE+IS10 validation row. Although paper-comparable cross-validation rows report fold-level variability, the core participant, temporal, and external-transfer ladder should ideally include DeLong, bootstrap, or paired-resampling confidence intervals. The large validation drops motivate caution, but small pairwise model differences should not be overinterpreted without those intervals.

The model-development bank included classical acoustic descriptors, boosted and kernel classifiers, fusion, CNN-style branches, self-supervised representations, and feature-search variants, but it did not exhaust every possible transformer, domain-adaptation, or event-segmentation design. Future gains in internal performance remain possible. The unresolved biomedical requirement is stronger evidence under temporally separated, externally matched, calibrated, and subgroup-aware validation. Calibration, device/country stratification, and prospective operating-point validation should be expanded before clinical use is considered.

The next-stage study should be designed around this validation gap. It should collect or identify an external cohort with matched cough, breath, and speech prompts, preserve device and site metadata, define COVID-19 labels through a known clinical process, predefine screening thresholds, and evaluate calibration and subgroup behavior before any deployment claim. Under that design, additional architectures such as spectrogram transformers or domain-adaptation models would be useful only if they improve performance at the temporal and external validation tiers, not only on a source-domain split. The figures in this manuscript are original study figures and do not reuse figures from prior publications.

\section{Conclusion}

COVID-19 respiratory-audio models can achieve strong internal discrimination while remaining fragile under the validation conditions that matter for screening. In this study, the selected Coswara fusion model reached AUROC 0.897, but performance declined under time-aware and strict temporal validation. In the fair cough-only comparison, Coswara internal AUROC of 0.849-0.868 contrasted with COUGHVID transfer AUROC of 0.523-0.543. Metadata-only models were highly predictive, with full safe metadata AUROC 0.964. These findings support a conservative conclusion: respiratory-audio screening research should be evaluated as a reliability and confounding problem, with explicit modality scope, temporal testing, external transfer, and metadata controls before clinical claims are made.

\section*{Data and Code Availability}

All numerical results were computed from the reproducible analysis pipeline. Public source datasets should be accessed through the original Coswara and COUGHVID maintainers under their respective terms. Derived aggregate tables, figure source data, manuscript figures, and analysis code are intended for release through the project repository or an archival public repository subject to dataset redistribution constraints. No new experimental results are introduced outside the reproducible analysis workflow.

\bibliographystyle{elsarticle-harv}
\bibliography{references}

\end{document}
